%!TEX root = ../../csuthesis_main.tex
\chapter{绪论}

\section{研究背景}

近年来，随着低空经济成为国家战略性新兴产业，无人机（Unmanned Aerial Vehicle，UAV）技术进入了快速发展阶段。2024 年，“低空经济”首次被写入政府工作报告，无人机作为低空领域最具代表性的智能装备，正从早期的航拍娱乐工具迅速渗透到电力巡检、应急救援、农业植保、地理测绘、城市安防等众多行业。随着机载相机、图传链路与飞行控制系统的持续进步，无人机能够在复杂环境中快速获取高分辨率航拍图像，为地面目标信息的低成本、高机动获取提供了新的技术手段。

在这一背景下，如何让无人机“看得懂”航拍图像，并进一步回答“目标在哪里”，成为无人机自主作业的关键。航拍图像目标检测负责从图像中定位并识别感兴趣目标，而地面定位负责将检测得到的像素位置换算为地面坐标，二者共同构成无人机“感知—定位—决策”链条的基础环节。然而，航拍图像具有成像高度大、视场范围广、目标尺度小且分布密集等特点，地面目标在不同飞行高度与成像视角下呈现出显著的形态差异，使得通用目标检测模型难以直接取得理想效果；同时，机载平台在算力、功耗与重量上的严格约束，又对检测模型的参数量与推理速度提出了更高要求。

此外，仅依靠单目航拍图像获取目标地面坐标，需要借助相机内参、飞行高度与成像几何关系建立像素到地面的映射模型，其精度会随飞行高度与成像视角的变化而显著衰减，而该类误差在实际作业中缺乏系统的量化分析。真实飞行试验成本高、周期长，且难以获得精确的目标真值，为上述问题的研究带来了困难。OpenHUTB 是面向智能体仿真与教学科研的开源平台，已集成无人车、无人机、具身人、水下机器人等多类载体与场景，可提供可控、可重复且带有仿真真值的实验环境，为航拍图像目标检测与定位方法的研究提供了理想的验证条件。基于此，本文以 OpenHUTB 无人机载体为依托，构建覆盖多高度、多视角的航拍数据集，研究轻量级目标检测与单目地面定位方法，并实现集检测显示、坐标输出与误差统计于一体的桌面系统。

\section{国内外研究现状}

\subsection{航拍图像目标检测研究现状}

目标检测是计算机视觉的基础任务之一，其发展大致经历了从手工特征到深度学习、从两阶段到单阶段、从卷积网络到 Transformer 的演进过程。Girshick 等提出的 R-CNN 系列将候选区域与卷积特征结合，奠定了两阶段检测器的基础；Ren 等提出的 Faster R-CNN 以区域建议网络实现端到端训练，显著提升了检测精度与速度。为满足实时性需求，Redmon 等提出 YOLO（You Only Look Once），将检测建模为统一的回归问题，实现了单阶段实时检测；Liu 等提出的 SSD 则通过多尺度特征图直接预测目标，进一步平衡了精度与速度。此后，YOLO 系列持续迭代，在锚框机制、标签分配、损失函数与网络结构等方面不断改进，其中 YOLOX 引入解耦头与无锚框设计，YOLOv7、YOLOv8 在训练策略与特征融合上进一步优化，YOLOv10 则从后处理环节入手实现端到端实时检测。与此同时，Carion 等提出的 DETR 将 Transformer 引入检测任务，摆脱了手工设计的后处理组件，RT-DETR 进一步将其改造为实时检测器，形成了 CNN 与 Transformer 两条并行技术路线。

航拍图像目标检测是上述通用方法在特定成像条件下的延伸。与自然场景图像相比，航拍图像由高空俯视或斜视成像获得，具有目标尺度小、分布密集、背景复杂、尺度跨度过大等特点，通用检测器直接迁移时往往出现漏检与误检。围绕这些问题，国内外研究主要从三个方向展开：一是多尺度特征增强，Lin 等提出的特征金字塔网络（FPN）及后续的 PANet、BiFPN 等结构，通过自顶向下与自底向上的路径融合深浅层特征，缓解小目标特征在多次下采样中丢失的问题；二是针对密集小目标的检测框架，如 ClusDet、DMNet 等通过区域拆分与密度估计提升密集场景的召回率，TPH-YOLOv5 则结合切片推理与注意力机制改善 VisDrone 等数据集上的小目标检测效果；三是推理阶段的工程改进，Akyon 等提出的 SAHI 采用切片辅助推理，将高分辨率航拍图切分为小块分别检测再融合，以较小代价提升小目标召回。

数据层面，航拍检测的公开基准不断丰富。Xia 等构建的 DOTA 面向遥感图像中的旋转目标检测，Lam 等发布的 xView 聚焦大幅面卫星图像中的小目标，Du 等建立的 UAVDT 与天津大学团队发起的 VisDrone 则以无人机视角的视频与图像为对象，成为无人机检测与跟踪研究中使用最广泛的基准之一；AI-TOD 等数据集则专门针对极小目标场景。这些数据集推动了航拍检测方法的快速发展，但其成像高度与视角相对固定，难以支撑“同一场景、不同高度与视角”的控制变量研究，且真值多为标注框而非目标的地理坐标，难以直接用于定位精度的评估

\subsection{单目视觉定位与像素—地面坐标换算研究现状}

将图像中的像素位置换算为地面坐标，是无人机执行巡检、搜救等任务时的共性需求。按照技术路线，现有方法可分为几何法、学习法与多传感器融合法三类。

几何法以相机成像模型为基础。Zhang 提出的棋盘格标定法建立了相机内参的标准化求解流程，使焦距、主点与畸变系数的获取不再依赖专用设备，成为后续几何定位工作的基础。在已知相机内参与飞行高度的前提下，若假设目标位于水平地面，则可通过相似三角形或单应性变换建立像素坐标与地面坐标的解析关系，逆透视变换（IPM）即是其中的典型代表，被广泛应用于车道线检测、交通监控与无人机目标定位。这类方法计算量小、可解释性强，但其精度高度依赖若干理想假设：相机光轴需与地面保持已知夹角、地面需近似为平面、飞行高度需准确可知。当无人机以斜视视角拍摄远距离目标时，上述假设被明显破坏，微小的姿态角误差会被放大为显著的地面定位偏差。

学习法试图绕开显式的几何建模。一类思路是使用单目深度估计网络直接回归场景深度，Godard 等提出的自监督深度估计方法与 Yang 等提出的 Depth Anything 系列模型，在通用场景上取得了良好的深度预测效果，为单目测距提供了新途径；另一类思路是以图像为输入、以目标位置为输出，端到端训练定位网络。然而，单目深度估计存在固有的尺度歧义，且模型在仿真数据与真实航拍数据之间存在明显的域偏移，直接用于米级精度的地面定位仍存在较大风险。

多传感器融合法则通过引入 GNSS、IMU、激光雷达或双目视觉，弥补单目图像缺乏深度信息的不足，能够在实际作业中取得较高精度，但随之带来成本、重量与功耗的上升，与中小型无人机对载荷的约束相矛盾。总体来看，在机载算力与载荷受限的条件下，基于已知飞行高度的单目几何定位仍是性价比最高的方案，但其误差随高度与视角变化的定量规律尚缺乏系统研究，这恰是本文关注的重点之一。

\subsection{轻量化检测模型与机载部署研究现状}

无人机机载平台的算力、功耗与散热条件有限，检测模型必须在精度与速度之间取得平衡。轻量化研究主要沿三条路径推进：一是轻量骨干网络设计，Howard 等提出的 MobileNet 系列以深度可分离卷积大幅降低计算量，Han 等提出的 GhostNet 通过廉价线性变换生成冗余特征图，ShuffleNet 则以通道混洗降低分组卷积的副作用；二是模型压缩，包括通道剪枝、量化感知训练与知识蒸馏，Hinton 等提出的知识蒸馏通过教师—学生结构将大模型的知识迁移至小模型，在保持精度的同时显著压缩模型体积；三是推理引擎优化，借助 TensorRT、ONNX Runtime、NCNN 等工具对算子融合、内存布局与量化策略进行针对性优化，使模型能够在 Jetson 等边缘设备上达到实时推理。

在航拍检测任务中，轻量化的代价往往体现为小目标精度的明显下降：浅层特征被过度压缩后，小目标在高层特征图上几乎消失。因此，近年来出现了以 YOLOv8n、YOLOv10n 为代表的“小模型 + 强数据增强”组合，以及将轻量骨干与多尺度注意力、特征重采样相结合的改进方案，试图在参数量与召回率之间寻找更优的平衡点。但现有工作大多在固定分辨率与固定视角的公开数据集上验证，缺少在统一硬件平台（如 RTX 4060 级显卡或同等算力的边缘设备）上对不同高度、不同视角条件下的精度与速度进行联合评估。

\subsection{研究述评}
综合上述三方面研究可以看出，航拍图像目标检测与地面定位已分别形成了较为成熟的技术体系，但仍存在以下不足：其一，数据层面缺乏覆盖多飞行高度、多成像视角且带有精确仿真真值的可控数据集，公开基准的成像条件相对单一，难以支撑尺度与视角对检测性能影响的定量分析；其二，方法层面检测与定位长期割裂研究，检测模型追求精度与速度，定位模型追求几何精度，二者缺乏统一的评估口径与一体化设计；其三，验证层面缺少以真值为基准、系统考察定位误差随高度与视角变化规律的工作，单目几何定位的适用边界尚不清晰；其四，工程层面多为算法验证脚本，缺少面向现场使用的桌面端一体化系统。

针对上述问题，本文以 OpenHUTB 无人机仿真平台为依托，构建覆盖三档高度、两类视角并带真值的航拍数据集，在统一硬件平台上研究轻量级检测模型与像素—地面坐标换算方法，定量分析检测精度与定位误差随高度、视角的变化规律，并基于 Qt 实现检测显示、坐标输出与误差统计一体化的桌面系统，以期在上述几个方面形成补充。

\section{研究意义}

研究一种面向机载部署的航拍图像目标检测与地面定位方法，对于提升无人机自主感知能力、降低作业成本、拓展低空应用场景具有重要的理论意义和现实意义

\subsection{理论意义}

基于 OpenHUTB 仿真平台的航拍图像目标检测与单目定位研究，将深度学习目标检测与几何视觉定位统一到同一可复现的实验框架之中，形成了“数据构建—模型检测—坐标换算—误差分析”的闭环研究范式。通过对低、中、高三档飞行高度与俯视、斜视两类视角的系统性控制变量实验，可以定量刻画目标尺度与成像视角对检测精度的影响规律，为航拍小目标检测的模型设计提供依据。

在定位环节，本文依据相机内参、飞行高度与目标像素坐标建立像素到地面坐标的换算模型，并分析定位误差随高度与视角的变化规律，有助于揭示单目视觉定位在航拍场景中的适用边界与误差传播机理，为同类单目几何定位方法提供可参考的误差评估框架。此外，对轻量化网络在机载算力约束下精度与速度权衡关系的探讨，以及对仿真数据训练、真值评估这一研究路径的实践，也可为仿真到真实的迁移研究提供方法借鉴。

\subsection{现实意义}

轻量化的航拍目标检测方法能够显著降低对高算力机载设备的依赖，使中小型无人机在普通边缘计算平台上即可完成实时检测，从而降低硬件成本与部署门槛。在电力线路、光伏电站、桥梁隧道等基础设施巡检中，自动化的航拍检测可以替代大量人工登高作业，在提高巡检效率的同时降低人员安全风险；在应急救援场景中，系统能够快速给出被困人员、车辆等目标的地面坐标，为指挥调度争取宝贵时间。

基于 Qt 实现的桌面系统将检测结果、目标坐标与误差统计集中呈现，便于现场操作人员直观使用，也为后续与指挥调度平台对接预留了接口。同时，本文构建的数据集、换算模型与软件工具均可在 OpenHUTB 开源生态中复用，对相关课程教学与后续研究具有一定的支撑作用。随着低空应用的不断拓展，这种集检测与定位于一体的轻量化方案在巡检、测绘、安防等场景中将发挥越来越重要的作用。

\section{研究内容与论文结构}

本文围绕 OpenHUTB 无人机载体的航拍图像目标检测与地面定位展开，主要研究内容包括航拍数据集的构建与标注、轻量级目标检测方法、像素—地面坐标换算与误差分析、检测定位系统的设计与实现四个方面。全文共分为七章，各章内容与上述研究内容相对应，具体安排如下。

第一章 绪论。本章介绍无人机航拍图像目标检测与地面定位的研究背景，阐明本课题在理论与现实两个层面的研究意义，并说明本文的主要研究内容与章节安排，为后续研究奠定基础。

第二章 相关技术与研究现状。本章综述无人机航拍目标检测、单目视觉定位以及轻量化网络三个方向的研究现状，介绍 YOLO 系列模型与 OpenHUTB 仿真平台的关键技术，分析现有方法在目标尺度变化、成像视角差异、机载部署与坐标映射等方面的局限性，从而明确本文拟解决的关键问题。

第三章 航拍数据集构建与标注。本章对应研究内容中的数据构建部分，介绍 OpenHUTB 无人机仿真场景的搭建方法与飞行参数设置，说明低、中、高三档飞行高度与俯视、斜视两类视角下的图像采集方案（采集图像不少于 800 帧，每档不少于 200 帧），制定目标边界框标注规范并完成不少于 500 例标注，最后给出数据集的统计特性与训练集、测试集的划分结果。

第四章 轻量级目标检测与单目定位方法。本章是本文的方法核心，对应轻量级目标检测方法与像素—地面坐标换算两项研究内容。首先介绍基线检测模型的结构与训练策略，从数据增强、特征融合与网络轻量化等方面阐述面向航拍小目标的改进设计；随后标定相机内参，结合飞行高度与目标像素坐标推导像素到地面坐标的换算模型，并给出误差传播与评估方法，为后续实验提供方法支撑。

第五章 系统设计与实现。本章对应系统实现这项研究内容，完成系统的需求分析与总体架构设计，划分图像输入、目标检测、坐标换算、结果展示与误差统计等功能模块，基于 Qt（C++ 或 Python 绑定）实现桌面端界面，并说明模块间的数据流与关键实现技术。

第六章 实验与结果分析。本章介绍实验环境、数据集配置与评价指标，展示不同高度档位与成像视角下检测精度与推理速度的对比结果，通过消融实验验证改进策略的有效性，检验 mAP@0.5 不低于 0.80、RTX 4060 上推理速度不低于 30 FPS 的指标达成情况；随后以仿真真值为基准统计定位误差，验证定位误差不超过 1.0 m，并分析误差随飞行高度与成像视角的变化规律，最后给出系统的功能测试与性能测试结果。

第七章 总结与展望。本章总结本文的主要工作与研究成果，归纳研究结论，指出数据集规模、仿真与真实场景差异等方面的不足，并对真实航拍数据验证、多目标跟踪与机载边缘部署等后续研究方向进行展望。



